% articulo-09-colores-enlaces.tex
% Colores, hipervínculos, URL e imágenes (\pandocbounded).
%
% ORDEN DE CARGA CRÍTICO:
%   xcolor → hyperref → xurl
%   hyperref debe cargarse después de la mayoría de paquetes
%   (tocloft, fancyhdr, titlesec, enumitem…) y xurl siempre después de hyperref.
%
% linkblue es el color de los hipervínculos externos y las citas bibliográficas.
% También lo usa estilo-español.lua directamente en los \hyperref que genera
% (\color{linkblue}): no renombrar sin actualizar también el filtro.
%
% Variables YAML (metadatos estándar de pandoc):
%   author, title, subject, keywords, lang

% ── Colores ───────────────────────────────────────────────────────────────────

\usepackage{xcolor}
\usepackage{etoolbox}
\definecolor{linkblue}{RGB}{0, 70, 150}

% ── Hipervínculos (hyperref) ──────────────────────────────────────────────────

\usepackage{hyperref}
\hypersetup{
  colorlinks    = true,
  linkcolor     = black,      % referencias cruzadas internas (incluye entradas del índice)
  urlcolor      = linkblue,   % URL
  citecolor     = linkblue,   % citas bibliográficas
  filecolor     = black,
  anchorcolor   = black,
  pdfauthor     = {$author$},
  pdftitle      = {$title$},
  pdfsubject    = {$subject$},
  pdfkeywords   = {$for(keywords)$$keywords$$sep$, $endfor$},
  pdflang       = {$if(lang)$$lang$$else$es-ES$endif$},
  bookmarksopen = true
}

% \textsc en bookmarks PDF:
% Con LuaLaTeX + fontspec, \textsc usa infraestructura Lua no disponible en el
% contexto restringido de \pdfstringdef. Tanto \MakeUppercase como \uppercase
% producen cadena vacía allí. Solución: pasar el argumento sin modificar.
% El bookmark mostrará el texto en minúsculas («xvi» en vez de «XVI») pero
% no quedará vacío. Para forzar mayúsculas en un caso concreto:
%   \texorpdfstring{\textsc{xvi}}{XVI}

\pdfstringdefDisableCommands{%
  \def\textsc#1{#1}%
  \def\scshape{}%
}

% ── URL (xurl) ────────────────────────────────────────────────────────────────
% Las URLs literales (<url>) usan la fuente principal del documento.
% \normalfont respeta la familia activa (serif o sans según tipografia-sans),
% a diferencia de \urlstyle{rm} que forzaba siempre la familia romana.

\usepackage{xurl}
\renewcommand{\UrlFont}{\normalfont}

% Corrección de \addcontentsline con hyperref: pandoc/citeproc emite
% \addcontentsline antes de \phantomsection en la bibliografía, lo que hace
% que el bookmark y la entrada del índice apunten a la página anterior.
% Se parchea \addcontentsline para que emita \phantomsection él mismo.
\makeatletter
\pretocmd{\addcontentsline}{\phantomsection}{}{}
\makeatother

% ── Remisiones internas en azul ───────────────────────────────────────────────
% Se activa linkcolor = linkblue globalmente y se fuerza negro solo en el
% índice y en las llamadas de nota al pie.
%
% El índice usa \contentsline; se parchea para forzar linkcolor=black localmente.
% Las llamadas de nota al pie usan \@makefnmark; se parchea para envolver
% la llamada con linkcolor=black.

\hypersetup{linkcolor = linkblue}

\makeatletter
\AtBeginDocument{%
  % Índice en negro.
  \let\Hy@org@contentsline\contentsline
  \renewcommand{\contentsline}[4]{%
    {\hypersetup{linkcolor=black}\Hy@org@contentsline{#1}{#2}{#3}{#4}}%
  }%
  % Llamadas de nota al pie en negro.
  \let\Hy@org@makefnmark\@makefnmark
  \renewcommand{\@makefnmark}{%
    {\hypersetup{linkcolor=black}\Hy@org@makefnmark}%
  }%
}
\makeatother

% ── \pandocbounded ────────────────────────────────────────────────────────────
% Limita imágenes a \linewidth × \textheight sin ampliarlas.

\makeatletter
\newsavebox{\pandoc@box}
\providecommand{\pandocbounded}[1]{%
  \sbox{\pandoc@box}{#1}%
  \Gscale@div{\@tempa}{\linewidth}{\wd\pandoc@box}%
  \Gscale@div{\@tempb}{\textheight}{\ht\pandoc@box}%
  \ifdim\@tempa pt > 1pt \def\@tempa{1}\fi%
  \ifdim\@tempb pt > 1pt \def\@tempb{1}\fi%
  \ifdim\@tempa pt >\@tempb pt%
    \scalebox{\@tempb}{\usebox{\pandoc@box}}%
  \else
    \scalebox{\@tempa}{\usebox{\pandoc@box}}%
  \fi%
}
\makeatother
